\section{Sieve and prime-distribution inputs}\label{sec:inputs}

We state the forms of the classical estimates used in the proof.
Appendix~\ref{sec:elementary-inputs} supplies elementary proofs of the
Buchstab normalization, the required additive large-sieve bound and the
plane-curve intersection estimate, and derives the progression upper
bound from the fundamental lemma. The inverse and stopping arguments
in the later sections are proved directly.

\begin{lemma}[Fundamental lemma of the sieve]\label{lem:fundamental}
Let $\mathcal P$ be any subset of the primes at most $v$. Suppose that
intersection counts for bad conditions indexed by squarefree products
$l$ of primes in $\mathcal P$ have the form
\[
 Xg(l)+r_l,
\]
where $g$ is multiplicative and
\[
 0\le g(p)\le\frac2p,\qquad g(p)<1,\qquad g(2)\le\frac12
 \quad\text{if }2\in\mathcal P.
\]
There are absolute constants $c,C>0$ and $s_0$ such that for $s\ge s_0$
the count avoiding all the bad conditions is
\begin{equation}\label{eq:fundamental-input}
 X\prod_{p\in\mathcal P}(1-g(p))\bigl(1+O(e^{-cs})\bigr)
 +O\!\left(
 \sum_{\substack{l\le v^s\\l\text{ squarefree}\\p\mid l\Rightarrow p\in\mathcal P}}
 \tau(l)^C|r_l|\right).
\end{equation}
The same statement holds for nonnegative weighted counts. A local
condition that makes the sifted set empty may instead be removed from
consideration as an exactly empty case.
\end{lemma}

This is a standard dimension-two form of the fundamental lemma; see
\cite[Lemma~2.8, p.~S137]{Sofos2023}, which records the stronger remainder sum
without divisor weights and attributes the result to
\cite[Corollary~6.10]{FriedlanderIwaniec2010}. To check the uniform
dimension condition here, put $g(p)=0$ outside $\mathcal P$. For $p\ge3$,
\[
 (1-2/p)^{-1}=(1-1/p)^{-2}\bigl(1+O(p^{-2})\bigr).
\]
Mertens' theorem therefore gives
\[
 \prod_{u\le p<v}(1-g(p))^{-1}
 \le K\left(\frac{\log v}{\log u}\right)^2
 \qquad(2\le u<v)
\]
with an absolute $K$. The factor at 2 is at most 2. The sieve weights
depend only on intersections of the bad conditions, so their use is
not restricted to literal divisibility conditions.

\begin{lemma}[Prime-distribution estimates]\label{lem:prime-inputs}
The following estimates will be used, with absolute constants unless
fixed subscripts indicate a dependence.
\begin{enumerate}[label=\textup{(\roman*)}]
\item There is $c>0$ such that
\[
 \pi(v)=\li(v)+O\!\left(\frac v{\log v}e^{-c\sqrt{\log v}}\right),
 \qquad
 \prod_{p\le v}(1-1/p)
 =\frac{e^{-\gamma}}{\log v}
   \bigl(1+O(e^{-c\sqrt{\log v}})\bigr).
\]
Decreasing $c$ absorbs harmless powers of $\log v$.
\item For every fixed $A,M>0$, uniformly for
$q\le(\log v)^A$ and $(a,q)=1$,
\[
 \pi(v;q,a)=\frac{\li(v)}{\varphi(q)}
             +O_{A,M}\!\left(\frac v{(\log v)^M}\right).
\]
The implied constants may be ineffective.
\item For $(a,q)=1$ and $H/q$ exceeding a suitable absolute constant,
\[
 \#\{x<p\le x+H:p\equiv a\pmod q\}
 \ll\frac{H}{\varphi(q)\log(H/q)}.
\]
\item Uniformly for $v^{9/10}\le H\le v$ as $v\to\infty$,
\[
 \#\{v-H<p\le v\}\ge(1-o(1))\frac H{\log v}.
\]
\end{enumerate}
\end{lemma}

Part (i) is the classical zero-free-region PNT together with its Mertens
product consequence; see \cite{IwaniecKowalski2004}, and the stronger
explicit PNT error in \cite{FioriKadiriSwidinsky2023}.
Parts (ii) and (iii) are Siegel--Walfisz and Brun--Titchmarsh; see the
explicit statements in \cite[p.~75]{GranvilleSoundararajan} and
\cite[Theorem~1.4.4, p.~39]{GranvilleSoundararajan2014}, respectively.
Lemma~\ref{lem:elementary-brun-titchmarsh} derives the order bound in
part (iii) from the stated sieve input and Mertens' formula.
Part (iv) follows from the short-interval
PNT of Huxley, for example in Heath-Brown's proof
\cite[Theorem, p.~1365]{HeathBrown1982}.
To see that a fixed-exponent version suffices,
tile an interval of length at least $v^{9/10}$ by intervals ending at $x$
of length $x^{4/5}$. Stop before the final incomplete tile or before
$x<v^{1/2}$. The lost total length is
$O(v^{4/5}+v^{1/2})=o(H)$. All remaining endpoints tend to infinity
uniformly, so the fixed-exponent theorem sums over the tiles. Prime
powers contribute a negligible amount, and $\log p\le\log v$ gives the
stated prime-count lower bound.

We also use the additive large sieve in the following standard form:
for distinct reduced fractions with denominators at most $Q$, and
coefficients supported on $J$ consecutive integers,
\begin{equation}\label{eq:large-sieve-input}
 \sum_{q\le Q}\ \sum_{\substack{a\bmod q\\(a,q)=1}}
 \left|\sum_j b_j e(aj/q)\right|^2
 \ll(J+Q^2)\sum_j|b_j|^2,
 \qquad e(t)=e^{2\pi it}.
\end{equation}
See \cite{MontgomeryVaughan1973} and \cite[Theorem~3]{VaughanLargeSieve}.
Lemma~\ref{lem:elementary-large-sieve} proves the order bound needed here.

For the complete Kloosterman sum
\[
 \Kl(h,k;N)=\sum_{\substack{u\bmod N\\(u,N)=1}}
 e\!\left(\frac{hu+k\bar u}{N}\right)
\]
we use
\begin{equation}\label{eq:kloosterman-input}
 |\Kl(h,k;N)|\ll \tau(N)(h,k,N)^{1/2}N^{1/2}.
\end{equation}
This holds for arbitrary positive moduli, including prime powers, and
for noncoprime or zero frequencies; see
\cite[Lemma~7.1]{Lichtman2022} and
\cite[Corollary~11.12, p.~280]{IwaniecKowalski2004}. The gcd factor in
\eqref{eq:kloosterman-input} will be retained in all completions.

The probability input is the nonarithmetic iid key renewal theorem:
if positive iid increments have finite positive mean $\mu$, renewal
measure $U$, and $H$ is directly Riemann integrable, then
\begin{equation}\label{eq:renewal-input}
 (U*H)(v)\longrightarrow\frac1\mu\int_\R H(u)\,du.
\end{equation}
We apply the usual half-line theorem after translating tests that vanish
on a left half-line; see \cite[Section~2.7, Theorem~35]{Serfozo2009}.
The regeneration, finite-mean, nonarithmetic and direct-integrability
hypotheses needed here are proved in Section~\ref{sec:functions}.
There we also use the classical Buchstab limit
$\omega_{\mathrm{Bu}}(s)\to e^{-\gamma}$, with its independent local
proof in Lemma~\ref{lem:elementary-buchstab}; see also
\cite[Section~28.9, p.~236]{MontgomeryVaughanIII}.

\subsection{Small-prime counts in a progression}

\begin{lemma}\label{lem:small-sieve}
Let one forbidden class be given at every prime $p\le u$, where
$u\to\infty$. Count the points in a progression segment whose coordinate
interval has real length $J$, and whose step is coprime to every prime
$p\le u$. Put $V_{\mathrm{sm}}(u)=\prod_{p\le u}(1-1/p)$.
For each fixed $\eta>0$, uniformly when $J\ge u^\eta$, its sifted count
$N$ satisfies
\begin{equation}\label{eq:small-upper}
 N\ll_\eta JV_{\mathrm{sm}}(u).
\end{equation}
There are absolute $c_1,C_1>0$ and a sufficiently large fixed $v_0$
such that, writing $v=\log_uJ$, uniformly for $v\ge v_0$,
\begin{equation}\label{eq:small-relative}
 \left|\frac{N}{JV_{\mathrm{sm}}(u)}-1\right|\le C_1e^{-c_1v}.
\end{equation}
All estimates are independent of the progression's initial point,
step and forbidden classes.
\end{lemma}

\begin{proof}
For any squarefree $l$ supported on the small primes, the simultaneous
bad conditions specify one class of the progression coordinate modulo
$l$. Their count is $J/l+O(1)$, with an absolute error independent of
the step and the classes. Apply Lemma~\ref{lem:fundamental} at level
$J^{1/2}$. Its sieve parameter is $v/2$, and the remainder sum is
\[
 \ll J^{1/2}(1+\log J)^{C_2}
\]
for an absolute fixed $C_2$. After division by $JV_{\mathrm{sm}}(u)$, this is
\[
 \ll J^{-1/2}(1+\log J)^{C_2}\log u.
\]
For all sufficiently large $u$ this is at most $e^{-c_1v}$ uniformly
for $v\ge v_0$, on decreasing $c_1$ and increasing $v_0$ if necessary.
The relative fundamental-lemma error has the same form. This proves
\eqref{eq:small-relative}.

For \eqref{eq:small-upper}, sieve only by primes at most $u^\theta$,
where $0<\theta<1$ is fixed sufficiently small in terms of $\eta$.
At level $J^{1/2}$ the sieve parameter is at least
$\eta/(2\theta)$ and can therefore be made larger than $s_0$.
The upper fundamental lemma gives
\[
 N\ll JV_{\mathrm{sm}}(u^\theta)+J^{1/2}(1+\log J)^{C_2}
 \ll_\eta JV_{\mathrm{sm}}(u),
\]
using Mertens and $J\ge u^\eta$. The same calculation covers either
endpoint convention and nonintegral $J$, since each intersection still
has an $O(1)$ endpoint error.
\end{proof}

We will repeatedly apply \eqref{eq:small-upper} to short difference
pieces created by freezing an interval length. If $\xi>0$ is fixed and
the coordinate length of such a piece is comparable to $\xi J$, with
$\xi J\ge u^\eta$, its cost is $O_\eta(\xi JV_{\mathrm{sm}}(u))$. The constant does
not depend on $\xi$; the threshold from which the length inequality
holds may depend on that fixed width. The same observation applies to
a specified residue subsequence of the piece.
